\section{结论}
结论为“保留”。三类证据方向一致：三池回归系数与平均Rank IC均为正，成本后五层收益在全部 A 股和中证500中严格单调。全部 A 股平均Rank IC为0.0566，多空年化13.00\%，Q5相对内部基准的算术年化超额为5.31个百分点。BP最适用于全部 A 股，在低流通市值域最强；两个指数股票池方向相同、强度较弱。BP可以作为后续多因子研究的估值类候选因子。

BP的效果具有时变性。Q5相对内部基准有两段较深的失效：2019-04见顶后回撤11.92\%、29个月后创新高，2012-11见顶后回撤11.34\%、低于前高24个观察月。样本末端处于新的失效期：Q5相对净值低于2024-08高点，滚动12个月超额自2025-06起连续为负，累计回归系数自2024-08后未创新高。行业层面，在2021-07至2025-12窗口内，房地产、钢铁等5个行业的Q5跑输各自基准。多空收益约59\%来自低BP端的落后；只做多时，可获得的主要是5.31个百分点的多头超额。

三类证据衡量的对象不同：WLS反映加权线性斜率，Rank IC反映全截面排序，等权高层组合只使用尾部股票并扣费。因此正IC、正回归与个别年份Q5的弱超额可以同时出现，D8、D9中高层IR的回落也说明高期望收益与低超额波动需要分开评价。

进入多因子合成时须注意共线风险。原始BP与对数总市值的逐月相关在$-$0.392至0.316之间，并有四段持续两年以上的同号期；行业间BP中位数相差近10倍（图A）。合成前应先检验BP与规模、行业以及EP等其他估值或质量因子的相关冗余，并在新的时间窗与可执行成本下确认其增量。

数据与成本口径有三点说明：历史PB与行业取自当前可得的历史版本，未作逐版本时点重建；停牌、退市与合并端点按已声明政策估值；成本采用教材只数换手，不含冲击成本与融券约束。多空均值的异方差与自相关稳健（HAC）检验、行业多重检验校正与样本外评估可作为后续工作。
